\section{Presupuesto óptico}
\label{sec:presupuesto}
% Consigna 8: sentido descendente (OLT -> ONT, 1490 nm) y ascendente (ONT -> OLT, 1310 nm),
% con las especificaciones de los equipos seleccionados. Incluir:
% - Potencias mínima y máxima de transmisión.
% - Sensibilidad y nivel máximo admisible de los receptores.
% - Atenuación de la fibra para cada longitud de onda.
% - Pérdidas de splitters, conectores y empalmes.
% - Margen de ingeniería (envejecimiento, reparaciones, incertidumbres).
% Verificar en ambos sentidos:
% 1. Recorrido de MAYOR pérdida -> potencia recibida suficiente.
% 2. Recorrido de MENOR pérdida -> no superar el nivel máximo del receptor.
% Identificar estos recorridos en el plano, presentar tablas y proponer correcciones si algo no cumple.
% Se puede reutilizar la metodología y el formato de tablas de la Parte 1 (Parte_1/parteB.tex).

\subsection{Parámetros de los equipos y de la ODN}
\label{sec:parametros_presupuesto}
\pendiente{potencias, sensibilidades y pérdidas unitarias por longitud de onda.}

\subsection{Sentido descendente (OLT $\rightarrow$ ONT)}
\label{sec:descendente}
\pendiente{recorridos de mayor y menor pérdida en sentido descendente.}

\subsection{Sentido ascendente (ONT $\rightarrow$ OLT)}
\label{sec:ascendente}
\pendiente{recorridos de mayor y menor pérdida en sentido ascendente.}
